% --- CORRECCIÓN IMPORTANTE EN LA PRIMERA LÍNEA ---
% Agregamos "xcolor=table" para que funcione \rowcolor sin errores
\documentclass[aspectratio=169, 10pt, xcolor=table]{beamer}

% --- Configuración del Tema (estilo profesional como el ejemplo) ---
\usetheme{Madrid}
\usecolortheme{dolphin}

% --- Paquetes necesarios ---
\usepackage[utf8]{inputenc}
\usepackage[spanish]{babel}
\usepackage{amsmath, amssymb}
\usepackage{booktabs}
\usepackage{graphicx}
\usepackage{listings}
\usepackage{tikz}
\usetikzlibrary{arrows.meta, positioning, shapes.geometric}

% --- Ruta de Imágenes (crea la carpeta "imagenes" y coloca ahí tus archivos) ---
\graphicspath{{./imagenes/}}

% --- Estilo para código Python ---
\definecolor{codegreen}{rgb}{0,0.6,0}
\definecolor{codegray}{rgb}{0.5,0.5,0.5}
\definecolor{codepurple}{rgb}{0.58,0,0.82}
\definecolor{backcolour}{rgb}{0.95,0.95,0.92}
\lstset{
    language=Python,
    backgroundcolor=\color{backcolour},
    commentstyle=\color{codegreen},
    keywordstyle=\color{blue},
    numberstyle=\tiny\color{codegray},
    stringstyle=\color{codepurple},
    basicstyle=\ttfamily\scriptsize,
    breaklines=true,
    frame=single,
    showstringspaces=false
}

% --- Pie de página personalizado ---
\makeatletter

\setbeamertemplate{footline}{
  \leavevmode%
  \hbox{%
  \begin{beamercolorbox}[wd=.333333\paperwidth,ht=2.25ex,dp=1ex,center]{author in head/foot}%
    \usebeamerfont{author in head/foot}\insertshortauthor
  \end{beamercolorbox}%
  \begin{beamercolorbox}[wd=.333333\paperwidth,ht=2.25ex,dp=1ex,center]{title in head/foot}%
    \usebeamerfont{title in head/foot}Sesión 5
  \end{beamercolorbox}%
  \begin{beamercolorbox}[wd=.333333\paperwidth,ht=2.25ex,dp=1ex,right]{date in head/foot}%
    \usebeamerfont{date in head/foot}NLP \hfill \insertframenumber/\inserttotalframenumber \hspace*{0.3cm}
  \end{beamercolorbox}}%
  \vskip0pt%
}

\makeatother

% --- Datos de la portada ---
\title[SESIÓN 5 NLP]{\textbf{Sesión 5}}
\subtitle{Modelos Clásicos para NLP: Naive Bayes, SVM, HMM, MaxEnt y CRF}
\author[Prof. C. Sánchez]{Carlos César Sánchez Coronel}
\institute{\textbf{NLP}}
\date{}

% Logo opcional (descomentar si tienes la imagen)
% \titlegraphic{\includegraphics[height=1.5cm]{logo.png}}

% --- Eliminar la barra de navegación (opcional) ---
\setbeamertemplate{navigation symbols}{}

% --- Índice al inicio de cada sección ---
\AtBeginSection[]{
  \begin{frame}
    \frametitle{Contenido}
    \tableofcontents[currentsection]
  \end{frame}
}

% ============================================================================
% INICIO DEL DOCUMENTO
% ============================================================================
\begin{document}

% --- Portada ---
\begin{frame}
    \titlepage
\end{frame}

% --- Índice general ---
\begin{frame}{Contenido}
    \tableofcontents
\end{frame}

% ============================================================================
% SECCIÓN 1: INTRODUCCIÓN
% ============================================================================
\section{Introducción}

\begin{frame}{¿Por qué estudiar modelos clásicos?}
    \begin{itemize}
        \item Aunque el deep learning domina la investigación, los modelos clásicos siguen siendo relevantes por:
        \begin{itemize}
            \item \textbf{Eficiencia con datos pequeños}: funcionan bien con miles de ejemplos.
            \item \textbf{Interpretabilidad}: sus decisiones son más fáciles de explicar y depurar.
            \item \textbf{Línea base obligatoria}: todo proyecto serio debe comparar contra ellos.
            \item \textbf{Producción en entornos restringidos}: bajo costo computacional.
        \end{itemize}
    \end{itemize}
    % Basado en introducción del PDF
\end{frame}

% ============================================================================
% SECCIÓN 2: NAIVE BAYES
% ============================================================================
\section{Naive Bayes para Clasificación de Texto}

\begin{frame}{Fundamento matemático}
    \begin{block}{Teorema de Bayes}
        \[
        P(c|d) = \frac{P(c) \prod_{i=1}^{n} P(w_i|c)}{P(d)}
        \]
        con supuesto de independencia condicional:
        \[
        P(\mathbf{w}|c) = \prod_i P(w_i|c)
        \]
    \end{block}
    \begin{itemize}
        \item Variantes: \textbf{MultinomialNB} (frecuencias), \textbf{BernoulliNB} (presencia/ausencia), \textbf{GaussianNB} (continuas).
    \end{itemize}
    % Basado en PDF y Pregunta 1
\end{frame}

\begin{frame}{Estimación de parámetros}
    \begin{block}{Suavizado Laplace}
        \[
        P(w|c) = \frac{\text{freq}(w,c) + \alpha}{\sum_{w'\in V} \text{freq}(w',c) + \alpha |V|}
        \]
        \begin{itemize}
            \item $\alpha = 1$ es común (add-1).
            \item Evita probabilidades cero para palabras no vistas.
        \end{itemize}
    \end{block}
    % Basado en PDF y Pregunta 2
\end{frame}

\begin{frame}{Ejemplo numérico}
    Corpus:
    \begin{itemize}
        \item d1: "el gato caza ratones" (animales)
        \item d2: "el perro persigue gatos" (animales)
        \item d3: "el coche acelera rápido" (vehículos)
    \end{itemize}
    Clasificar "gato rápido" usando MultinomialNB con $\alpha=1$.
    \vspace{0.3cm}
    \begin{columns}[t]
        \column{0.5\textwidth}
        \textbf{Animales}
        \[
        P(\text{animal}) = \frac{2}{3}
        \]
        \[
        P(\text{gato}|\text{animal}) = \frac{1+1}{6+10} = 0.125
        \]
        \[
        P(\text{rápido}|\text{animal}) = \frac{0+1}{6+10} = 0.0625
        \]
        \[
        \text{score} \propto 0.00521
        \]
        \column{0.5\textwidth}
        \textbf{Vehículos}
        \[
        P(\text{vehículo}) = \frac{1}{3}
        \]
        \[
        P(\text{gato}|\text{vehículo}) = \frac{0+1}{4+10} = 0.0714
        \]
        \[
        P(\text{rápido}|\text{vehículo}) = \frac{1+1}{4+10} = 0.1429
        \]
        \[
        \text{score} \propto 0.00340
        \]
    \end{columns}
    Clasificación: \textbf{animales}.
    % Basado en PDF
\end{frame}

\begin{frame}{Ventajas y limitaciones}
    \begin{columns}[t]
        \column{0.5\textwidth}
        \textbf{Ventajas}
        \begin{itemize}
            \item Extremadamente rápido.
            \item Funciona bien con datos pequeños.
            \item Maneja alta dimensionalidad.
            \item Interpretable.
        \end{itemize}
        \column{0.5\textwidth}
        \textbf{Limitaciones}
        \begin{itemize}
            \item Asume independencia (violada en lenguaje).
            \item No captura relaciones entre palabras.
        \end{itemize}
    \end{columns}
    % Basado en PDF
\end{frame}

% ============================================================================
% SECCIÓN 3: SUPPORT VECTOR MACHINES (SVM)
% ============================================================================
\section{Support Vector Machines (SVM) para Texto}

\begin{frame}{SVM lineal para texto}
    \begin{block}{Formulación}
        \[
        \min_{\mathbf{w},b} \frac{1}{2}\|\mathbf{w}\|^2 + C\sum_i \xi_i
        \]
        sujeto a $y_i(\mathbf{w}^T\mathbf{x}_i + b) \ge 1 - \xi_i$, $\xi_i \ge 0$.
    \end{block}
    \begin{itemize}
        \item $C$ controla penalización por errores.
        \item Para texto, kernel lineal suele ser suficiente por alta dimensionalidad.
    \end{itemize}
    % Basado en PDF y Pregunta 1, 9
\end{frame}

\begin{frame}{¿Por qué SVM funciona bien en texto?}
    \begin{itemize}
        \item \textbf{Alta dimensionalidad}: los textos tienen miles de características, SVM maneja $p \gg n$.
        \item \textbf{Representaciones dispersas}: BOW/TF-IDF son esparsas, SVM las maneja eficientemente.
        \item \textbf{Margen máximo}: buena generalización.
        \item \textbf{Kernel lineal suficiente}: los datos de texto suelen ser aproximadamente separables linealmente.
    \end{itemize}
\end{frame}

\begin{frame}[fragile]{Pipeline típico con scikit-learn}
    \begin{lstlisting}
from sklearn.feature_extraction.text import TfidfVectorizer
from sklearn.svm import LinearSVC
from sklearn.pipeline import make_pipeline

docs = ["el gato caza ratones", "el perro persigue gatos", "el coche acelera rápido"]
y = ["animales", "animales", "vehículos"]

pipeline = make_pipeline(
    TfidfVectorizer(),
    LinearSVC(C=1.0, class_weight='balanced')
)
pipeline.fit(docs, y)
    \end{lstlisting}
\end{frame}

\begin{frame}{Ventajas y limitaciones de SVM}
    \begin{columns}[t]
        \column{0.5\textwidth}
        \textbf{Ventajas}
        \begin{itemize}
            \item Robusto, buen rendimiento.
            \item Maneja desbalance con pesos.
            \item Bueno para datos de alta dimensión.
        \end{itemize}
        \column{0.5\textwidth}
        \textbf{Limitaciones}
        \begin{itemize}
            \item Menos interpretable que Naive Bayes.
            \item Requiere normalización de características.
        \end{itemize}
    \end{columns}
\end{frame}

% ============================================================================
% SECCIÓN 4: MODELOS OCULTOS DE MARKOV (HMM)
% ============================================================================
\section{Modelos Ocultos de Markov (HMM)}

\begin{frame}{Definición formal de HMM}
    \begin{block}{Componentes}
        \begin{itemize}
            \item Estados ocultos $Q = \{q_1,\dots,q_N\}$ (ej. etiquetas POS).
            \item Observaciones $O = \{o_1,\dots,o_T\}$ (ej. palabras).
            \item Matriz de transición $A = [a_{ij}]$ con $a_{ij} = P(q_t=j|q_{t-1}=i)$.
            \item Matriz de emisión $B = [b_j(k)]$ con $b_j(k) = P(o_t=k|q_t=j)$.
            \item Distribución inicial $\pi = [\pi_i]$ con $\pi_i = P(q_1=i)$.
        \end{itemize}
    \end{block}
    % Basado en PDF y Pregunta 3
\end{frame}

\begin{frame}{La propiedad de Markov}
    \begin{block}{Propiedad de Markov de primer orden}
        \[
        P(q_{t+1}|q_t, q_{t-1}, \dots, q_1) = P(q_{t+1}|q_t)
        \]
        El estado futuro solo depende del estado actual.
    \end{block}
    \begin{itemize}
        \item Limitación: no considera contexto más amplio.
    \end{itemize}
\end{frame}

\begin{frame}{Tres problemas fundamentales en HMM}
    \begin{table}
        \centering
        \begin{tabular}{l|c}
            \toprule
            \textbf{Problema} & \textbf{Algoritmo} \\
            \midrule
            Evaluación: $P(O|\lambda)$ & Forward \\
            Decodificación: mejor secuencia de estados & Viterbi \\
            Aprendizaje: estimar $\lambda = (A,B,\pi)$ & Baum-Welch (EM) \\
            \bottomrule
        \end{tabular}
    \end{table}
    % Basado en PDF y Pregunta 8
\end{frame}

\begin{frame}[fragile]{POS tagging con HMM en NLTK}
    \begin{lstlisting}
import nltk
from nltk.tag import HiddenMarkovModelTrainer
from nltk.corpus import treebank

train_data = treebank.tagged_sents()[:3900]
tagger = HiddenMarkovModelTrainer().train_supervised(train_data)

sentence = "The cat sits on the mat".split()
tags = tagger.tag(sentence)
print(tags)  # [('The', 'DT'), ('cat', 'NN'), ...]
    \end{lstlisting}
    % Basado en PDF
\end{frame}

\begin{frame}{Aplicaciones y limitaciones de HMM}
    \begin{itemize}
        \item \textbf{Aplicaciones}: POS tagging, NER, reconocimiento de voz, segmentación.
        \item \textbf{Ventajas}: modelo generativo, interpretable, funciona con datos etiquetados.
        \item \textbf{Limitaciones}: asume independencia de observaciones dado el estado, dependencias de primer orden.
    \end{itemize}
    % Basado en PDF y Pregunta 4
\end{frame}

% ============================================================================
% SECCIÓN 5: MODELOS DE MÁXIMA ENTROPÍA (MAXENT / REGRESIÓN LOGÍSTICA)
% ============================================================================
\section{Modelos de Máxima Entropía (MaxEnt)}

\begin{frame}{Fundamento matemático}
    \begin{block}{Probabilidad condicional}
        \[
        P(c|d) = \frac{\exp\left(\sum_j \lambda_j f_j(c,d)\right)}{\sum_{c'}\exp\left(\sum_j \lambda_j f_j(c',d)\right)}
        \]
        donde $f_j$ son funciones características arbitrarias.
    \end{block}
    \begin{itemize}
        \item Equivalente a regresión logística multinomial.
        \item No asume independencia; puede incluir características correlacionadas.
    \end{itemize}
    % Basado en PDF y Pregunta 7
\end{frame}

\begin{frame}{Características (features)}
    \begin{itemize}
        \item Pueden ser arbitrarias: unigramas, bigramas, contexto, sufijos, etc.
        \item Ejemplo: $f(\text{positivo}, d) = 1$ si la palabra "excelente" aparece en $d$.
        \item Permite modelar dependencias complejas sin supuestos de independencia.
    \end{itemize}
\end{frame}

\begin{frame}{Entrenamiento y regularización}
    \begin{block}{Log-verosimilitud negativa con regularización}
        \[
        \mathcal{L} = -\sum_i \log P(c_i|d_i) + \frac{1}{2C}\sum_j \lambda_j^2
        \]
    \end{block}
    \begin{itemize}
        \item Optimización con gradiente descendente o L-BFGS.
        \item Regularización L2 (Ridge) o L1 (Lasso) evitan overfitting.
    \end{itemize}
    % Basado en Pregunta 17
\end{frame}

\begin{frame}[fragile]{Ejemplo con scikit-learn}
    \begin{lstlisting}
from sklearn.linear_model import LogisticRegression
from sklearn.feature_extraction.text import TfidfVectorizer

X = ["buena película", "mala película", "excelente actuación"]
y = ["positivo", "negativo", "positivo"]

vectorizer = TfidfVectorizer()
X_vec = vectorizer.fit_transform(X)
clf = LogisticRegression(C=1.0, penalty='l2')
clf.fit(X_vec, y)
    \end{lstlisting}
\end{frame}

\begin{frame}{Ventajas y limitaciones de MaxEnt}
    \begin{itemize}
        \item \textbf{Ventajas}: modelo discriminativo, flexible, no asume independencia.
        \item \textbf{Limitaciones}: no modela secuencias explícitamente (ignora dependencias entre etiquetas consecutivas).
    \end{itemize}
\end{frame}

% ============================================================================
% SECCIÓN 6: CONDITIONAL RANDOM FIELDS (CRF)
% ============================================================================
\section{Conditional Random Fields (CRF)}

\begin{frame}{De HMM a CRF}
    \begin{itemize}
        \item HMM: generativo, modela $P(X,Y)$, asume independencia de observaciones.
        \item CRF: discriminativo, modela $P(Y|X)$ directamente, permite características arbitrarias y dependencias contextuales.
    \end{itemize}
\end{frame}

\begin{frame}{Definición formal de CRF}
    \[
    P(\mathbf{Y}|\mathbf{X}) = \frac{1}{Z(\mathbf{X})} \exp\left( \sum_{t=1}^{T} \sum_{k} \lambda_k f_k(y_t, y_{t-1}, \mathbf{X}, t) \right)
    \]
    \begin{itemize}
        \item $f_k$: funciones características (pueden depender de toda la secuencia de entrada).
        \item $\lambda_k$: pesos aprendidos.
        \item $Z(\mathbf{X})$: función de partición (normalización).
    \end{itemize}
    % Basado en PDF y Pregunta 13
\end{frame}

\begin{frame}{Tipos de características}
    \begin{itemize}
        \item \textbf{De estado}: relación etiqueta actual con observaciones (ej. "palabra termina en -ción").
        \item \textbf{De transición}: relación etiqueta actual con anterior (ej. "es poco probable que DT siga a JJ").
    \end{itemize}
\end{frame}

\begin{frame}[fragile]{Implementación con sklearn-crfsuite}
    \begin{lstlisting}
import sklearn_crfsuite
from sklearn_crfsuite import metrics

def word2features(sent, i):
    word = sent[i][0]
    return {
        'word': word,
        'is_capitalized': word[0].isupper(),
        'suffix-2': word[-2:],
        'prev_word': '' if i==0 else sent[i-1][0],
    }

X_train = [word2features(s, i) for s in train_sents ...]
crf = sklearn_crfsuite.CRF(algorithm='lbfgs', c1=0.1, c2=0.1)
crf.fit(X_train, y_train)
    \end{lstlisting}
    % Basado en PDF
\end{frame}

\begin{frame}{Ventajas de CRF sobre HMM}
    \begin{itemize}
        \item No asume independencia de las observaciones.
        \item Puede incluir características arbitrarias y superpuestas.
        \item Mejor rendimiento en tareas como NER y chunking.
        \item Útil cuando se dispone de suficientes datos y características ricas.
    \end{itemize}
    % Basado en Pregunta 4
\end{frame}

% ============================================================================
% SECCIÓN 7: COMPARACIÓN DE MODELOS
% ============================================================================
\section{Comparación de Modelos Clásicos}

\begin{frame}{Tabla comparativa}
    \begin{table}
        \centering
        \begin{tabular}{l|c|c|c}
            \toprule
            \textbf{Modelo} & \textbf{Tipo} & \textbf{Velocidad} & \textbf{Interpretabilidad} \\
            \midrule
            Naive Bayes & Generativo & Muy rápida & Alta \\
            SVM (lineal) & Discriminativo & Rápida & Media \\
            HMM & Generativo secuencial & Media & Alta \\
            MaxEnt & Discriminativo & Media & Media \\
            CRF & Discriminativo secuencial & Lenta & Baja \\
            \bottomrule
        \end{tabular}
    \end{table}
    % Basado en PDF
\end{frame}

\begin{frame}{¿Cuándo usar cada uno?}
    \begin{itemize}
        \item \textbf{Clasificación simple}: Naive Bayes (baseline), SVM (mejor rendimiento).
        \item \textbf{Etiquetado secuencial}: HMM (pocos datos), CRF (características ricas).
        \item \textbf{Cuando la independencia es insostenible}: MaxEnt o CRF.
    \end{itemize}
    % Basado en PDF
\end{frame}

% ============================================================================
% SECCIÓN 8: CASO INTEGRADO
% ============================================================================
\section{Caso Integrado: Clasificación y NER en Reseñas}

\begin{frame}{Problema}
    \begin{itemize}
        \item Corpus de reseñas de productos con calificación (1-5 estrellas).
        \item Tareas:
        \begin{enumerate}
            \item Clasificación de sentimiento (positivo/negativo).
            \item Extracción de aspectos mencionados (NER: batería, pantalla, precio).
        \end{enumerate}
    \end{itemize}
    % Basado en PDF
\end{frame}

\begin{frame}{Solución con modelos clásicos}
    \begin{columns}[t]
        \column{0.5\textwidth}
        \textbf{Clasificación}
        \begin{itemize}
            \item SVM con TF-IDF.
            \item MaxEnt con características adicionales.
            \item Comparar con Naive Bayes.
        \end{itemize}
        \column{0.5\textwidth}
        \textbf{NER}
        \begin{itemize}
            \item HMM simple.
            \item CRF con características ricas (contexto, sufijos).
        \end{itemize}
    \end{columns}
    \vspace{0.3cm}
    \textbf{Resultados esperados}: SVM supera a NB; CRF supera a HMM.
\end{frame}

% ============================================================================
% SECCIÓN 9: PREPARACIÓN PARA ENTREVISTAS
% ============================================================================
\section{Preparación para Entrevistas}

\begin{frame}{Preguntas típicas}
    \begin{itemize}
        \item ¿Cuándo usar Naive Bayes vs SVM? (datos pequeños vs. robustez)
        \item ¿Por qué SVM con kernel lineal funciona bien en texto?
        \item Explica el algoritmo de Viterbi.
        \item Diferencia entre HMM y CRF.
        \item ¿Qué es el suavizado Laplace?
        \item ¿Cómo manejarías desbalance en SVM?
    \end{itemize}
    % Basado en PDF y solucionario
\end{frame}

% ============================================================================
% SECCIÓN 10: CONEXIÓN CON EL RESTO DEL CURSO
% ============================================================================
\section{Conexión con el resto del curso}

\begin{frame}{Próximos pasos}
    \begin{itemize}
        \item \textbf{Sesión 4} (Word Embeddings): se pueden usar como características densas en estos modelos.
        \item \textbf{Sesión 6} (RNN/LSTM): superan a HMM/CRF en secuencias, pero requieren más datos.
        \item \textbf{Sesión 7} (Transformers): evolución natural para secuencias.
    \end{itemize}
\end{frame}

% ============================================================================
% SECCIÓN 11: RESUMEN
% ============================================================================
\section{Resumen}

\begin{frame}{Lo que hemos aprendido}
    \begin{exampleblock}{Conceptos clave}
        \begin{itemize}
            \item \textbf{Naive Bayes}: clasificación rápida, asume independencia.
            \item \textbf{SVM}: clasificación robusta, margen máximo, kernel lineal.
            \item \textbf{HMM}: modelado de secuencias generativo, Viterbi, Baum-Welch.
            \item \textbf{MaxEnt}: regresión logística multinomial, características arbitrarias.
            \item \textbf{CRF}: discriminativo para secuencias, supera a HMM en rendimiento.
            \item \textbf{Comparación}: elección según tarea, datos y necesidad de interpretabilidad.
            \item \textbf{Caso práctico}: clasificación y NER con modelos clásicos.
        \end{itemize}
    \end{exampleblock}
\end{frame}

% --- Diapositiva final ---
\begin{frame}
    \centering
    \Huge \textbf{¡Gracias!}

\end{frame}

\end{document}